\documentclass[11pt]{article}
\usepackage[margin=1in]{geometry}
\usepackage{amsmath,amssymb,mathtools,amsthm}
\usepackage{booktabs,longtable}
\usepackage{enumitem}
\usepackage[hidelinks]{hyperref}

\newcommand{\dd}{\mathrm{d}}
\newcommand{\R}{\mathbb{R}}
\newcommand{\TT}{\mathrm{TT}}
\newcommand{\E}{\mathbb{E}}
\newcommand{\calX}{\mathcal{X}}
\newcommand{\calU}{\mathcal{U}}
\newcommand{\calB}{\mathcal{B}}
\newcommand{\calD}{\mathcal{D}}
\newcommand{\dotp}{\partial_i}

% This teaching note cites many displayed formulas.  Number every display so
% cross-references point to a visible formula rather than to the surrounding
% section number.
\renewcommand{\[}{\begin{equation}}
\renewcommand{\]}{\end{equation}}

\newtheorem{proposition}{Proposition}
\newtheorem{lemma}{Lemma}
\newtheorem{corollary}{Corollary}
\theoremstyle{definition}
\newtheorem{definition}{Definition}
\newtheorem{assumption}{Assumption}
\theoremstyle{remark}
\newtheorem{remark}{Remark}

\title{Teaching the Fixed-Branch Zhao--Cui Tensor-Train Filter}
\author{BayesFilter P14 Pedagogical Mathematical Note}
\date{2026-05-31}

\begin{document}
\maketitle

\begin{abstract}
This note retells the fixed-branch Zhao--Cui tensor-train filtering
construction as a mathematical teaching document.  It is written for a
scientific reader who is comfortable with probability densities, Gaussian
likelihoods, integrals, least squares, and matrix factorizations, but who has
not previously worked with tensor-train filtering.  The presentation keeps the
P12--P13 mathematical content: the fixed-branch squared-TT recursion produces
normalized approximate filters, and its analytical gradient differentiates the
same approximate likelihood scalar computed by the forward pass.  The
difference is pedagogical.  Tensor trains are introduced first as low-rank
coupling in coefficient arrays; the filtering construction is derived before
the source papers are cited; and the gradient proof is split into layers so
that each derivative has a visible role in the filtering algorithm.
\end{abstract}

\section{Filtering Objects Before Tensor Notation}

The mathematical object of interest is not a tensor train.  It is the
unnormalized density that appears in one Bayesian filtering step.  Let
\[
  x_t\in\R^{d_x},\qquad \vartheta\in\R^{d_\vartheta},\qquad y_t\in\R^{d_y}
\]
denote the current state, a static parameter, and the observation.  A dominated
state-space model supplies
\[
  p_0(x_0,\vartheta),\qquad
  f_\alpha(x_t\mid x_{t-1},\vartheta),\qquad
  g_\alpha(y_t\mid x_t,\vartheta),
  \label{eq:p14-ssm}
\]
where \(\alpha\) is the external parameter with respect to which we may later
differentiate a likelihood.

Assume that, after observing \(y_{1:t-1}\), the filter has density
\[
  p_{t-1}(x_{t-1},\vartheta)
  =
  p(x_{t-1},\vartheta\mid y_{1:t-1}).
\]
Prediction and update can be written as
\[
  p(x_t,\vartheta\mid y_{1:t-1})
  =
  \int f_\alpha(x_t\mid x_{t-1},\vartheta)
       p_{t-1}(x_{t-1},\vartheta)\,\dd x_{t-1},
  \label{eq:p14-prediction}
\]
\[
  p_t(x_t,\vartheta)
  =
  \frac{
  g_\alpha(y_t\mid x_t,\vartheta)
  p(x_t,\vartheta\mid y_{1:t-1})}
  {p_\alpha(y_t\mid y_{1:t-1})}.
  \label{eq:p14-update}
\]
Now introduce the enlarged variable
\[
  u_t=(x_t,\vartheta,x_{t-1})
\]
and define
\[
  q_t(u_t;\alpha)
  =
  p_{t-1}(x_{t-1},\vartheta)
  f_\alpha(x_t\mid x_{t-1},\vartheta)
  g_\alpha(y_t\mid x_t,\vartheta).
  \label{eq:p14-q}
\]
Its integral is
\[
  Z_t(\alpha)
  =
  \int q_t(x_t,\vartheta,x_{t-1};\alpha)
  \,\dd x_t\,\dd\vartheta\,\dd x_{t-1}.
  \label{eq:p14-Z}
\]
Substituting \eqref{eq:p14-q} and integrating over \(x_{t-1}\) first gives
\[
\begin{aligned}
  Z_t(\alpha)
  &=
  \int g_\alpha(y_t\mid x_t,\vartheta)
  \left[
  \int f_\alpha(x_t\mid x_{t-1},\vartheta)
  p_{t-1}(x_{t-1},\vartheta)\,\dd x_{t-1}
  \right]
  \dd x_t\,\dd\vartheta \\
  &=
  \int g_\alpha(y_t\mid x_t,\vartheta)
  p(x_t,\vartheta\mid y_{1:t-1})\,\dd x_t\,\dd\vartheta
  =
  p_\alpha(y_t\mid y_{1:t-1}).
\end{aligned}
\]
The updated filter is therefore
\[
  p_t(x_t,\vartheta)
  =
  \int
  \frac{q_t(x_t,\vartheta,x_{t-1};\alpha)}
       {Z_t(\alpha)}
  \,\dd x_{t-1}.
  \label{eq:p14-filter-as-marginal}
\]

\paragraph{Why this matters.}
Every approximation in this note is an approximation to the three operations
in \eqref{eq:p14-Z}--\eqref{eq:p14-filter-as-marginal}: represent \(q_t\),
integrate it, and marginalize the old state.  Tensor trains are useful only if
they make those three operations feasible in high dimension.

\paragraph{Failure mode.}
If the approximation to \(q_t\) is poor in regions that contribute meaningful
mass to the integral, the output can be normalized but inaccurate.  Normalized
does not mean close to the exact posterior.

\section{Scalar Nonlinear Filtering Example}

Before using tensor notation, consider the scalar nonlinear model
\[
  x_t=\rho x_{t-1}+\eta_t,\qquad
  \eta_t\sim N(0,\sigma_\eta^2),
\]
\[
  y_t=x_t^2+\epsilon_t,\qquad
  \epsilon_t\sim N(0,\sigma_\epsilon^2).
  \label{eq:p14-scalar-model}
\]
Given a previous approximate filter \(\widehat p_{t-1}(x_{t-1})\), the
one-step unnormalized density is
\[
  q_t(x_t,x_{t-1})
  =
  \widehat p_{t-1}(x_{t-1})
  \frac{1}{\sigma_\eta}
  \varphi\!\left(\frac{x_t-\rho x_{t-1}}{\sigma_\eta}\right)
  \frac{1}{\sigma_\epsilon}
  \varphi\!\left(\frac{y_t-x_t^2}{\sigma_\epsilon}\right),
  \label{eq:p14-scalar-q}
\]
where \(\varphi\) is the standard normal density.

If \(y_t>0\) and \(\sigma_\epsilon\) is small, then
\[
  \frac{1}{\sigma_\epsilon}
  \varphi\!\left(\frac{y_t-x_t^2}{\sigma_\epsilon}\right)
\]
has large values near both \(x_t=\sqrt{y_t}\) and
\(x_t=-\sqrt{y_t}\).  Thus a single observation can create two plausible
state regions.  The non-Gaussian shape is visible directly in the formula; it
does not require any tensor argument.

Let \(\phi_t\) approximate the square root:
\[
  \phi_t(x_t,x_{t-1})\approx \sqrt{q_t(x_t,x_{t-1})}.
  \label{eq:p14-scalar-root}
\]
Define the nonnegative approximation
\[
  \widehat q_t(x_t,x_{t-1})
  =
  \phi_t(x_t,x_{t-1})^2
  +
  \tau_t\lambda_x(x_t)\lambda_{x^-}(x_{t-1}),
  \label{eq:p14-scalar-qhat}
\]
where
\[
  \tau_t\ge0,\qquad
  \int\lambda_x(x_t)\,\dd x_t=1,\qquad
  \int\lambda_{x^-}(x_{t-1})\,\dd x_{t-1}=1.
\]
The approximate normalizer is
\[
\begin{aligned}
  \widehat Z_t
  &=
  \iint \widehat q_t(x_t,x_{t-1})
  \,\dd x_t\,\dd x_{t-1}  \\
  &=
  \iint \phi_t(x_t,x_{t-1})^2\,\dd x_t\,\dd x_{t-1}
  +
  \tau_t
  \left[\int\lambda_x(x_t)\,\dd x_t\right]
  \left[\int\lambda_{x^-}(x_{t-1})\,\dd x_{t-1}\right] \\
  &=
  \iint \phi_t(x_t,x_{t-1})^2\,\dd x_t\,\dd x_{t-1}
  +\tau_t.
\end{aligned}
  \label{eq:p14-scalar-Zhat}
\]
The next approximate filter is
\[
\begin{aligned}
  \widehat p_t(x_t)
  &=
  \int
  \frac{\widehat q_t(x_t,x_{t-1})}{\widehat Z_t}
  \,\dd x_{t-1} \\
  &=
  \frac{1}{\widehat Z_t}
  \int\phi_t(x_t,x_{t-1})^2\,\dd x_{t-1}
  +
  \frac{\tau_t}{\widehat Z_t}\lambda_x(x_t).
\end{aligned}
  \label{eq:p14-scalar-next-filter}
\]

\paragraph{Why this matters.}
The defensive term is not mysterious.  It contributes a controlled background
density \((\tau_t/\widehat Z_t)\lambda_x(x_t)\) to the next filter.  The main
shape still comes from the square-root fit through
\(\int\phi_t^2\,\dd x_{t-1}\).

\paragraph{Failure mode.}
If \(\tau_t\) is too large, the defensive density can dominate the filter.  If
\(\tau_t\) is too small and \(\phi_t\) misses an important region, the filter
can assign too little mass there.  The equations show the tradeoff explicitly.

\section{Tensor Trains As Low-Rank Coupling, Not Magic Compression}

Start with two variables \(u_1,u_2\) and basis functions
\(\{b_{1,a}\}_{a=1}^{n_1}\), \(\{b_{2,b}\}_{b=1}^{n_2}\).  A full coefficient
expansion is
\[
  \phi(u_1,u_2)
  =
  \sum_{a=1}^{n_1}\sum_{b=1}^{n_2}
  F_{ab} b_{1,a}(u_1)b_{2,b}(u_2).
  \label{eq:p14-two-full}
\]
The coefficient matrix \(F\in\R^{n_1\times n_2}\) stores how coordinate
\(u_1\) couples to coordinate \(u_2\).  A rank-\(r\) factorization replaces
\[
  F_{ab}
  \approx
  \sum_{\ell=1}^r U_{a\ell}V_{\ell b}.
  \label{eq:p14-two-rank}
\]
Substituting into \eqref{eq:p14-two-full} gives
\[
\begin{aligned}
  \phi(u_1,u_2)
  &\approx
  \sum_{\ell=1}^r
  \left[\sum_{a=1}^{n_1}U_{a\ell}b_{1,a}(u_1)\right]
  \left[\sum_{b=1}^{n_2}V_{\ell b}b_{2,b}(u_2)\right] \\
  &=
  \sum_{\ell=1}^r U_\ell(u_1)V_\ell(u_2).
\end{aligned}
  \label{eq:p14-two-functional-rank}
\]
The storage changes from \(n_1n_2\) numbers to \(r(n_1+n_2)\) numbers.

For three variables, a full coefficient tensor
\[
  F_{abc},\qquad
  1\le a\le n_1,\quad 1\le b\le n_2,\quad 1\le c\le n_3,
\]
can be approximated by the chain
\[
  F_{abc}
  \approx
  \sum_{\ell_1=1}^{r_1}\sum_{\ell_2=1}^{r_2}
  G_1(a,\ell_1)G_2(\ell_1,b,\ell_2)G_3(\ell_2,c).
  \label{eq:p14-three-tt}
\]
The corresponding function is
\[
  \phi(u_1,u_2,u_3)
  \approx
  G_1(u_1)G_2(u_2)G_3(u_3),
  \label{eq:p14-three-functional}
\]
where \(G_1(u_1)\) is a \(1\times r_1\) row, \(G_2(u_2)\) is an
\(r_1\times r_2\) matrix, and \(G_3(u_3)\) is an \(r_2\times1\) column.
Multiplication produces a scalar.

The \(D\)-coordinate tensor train is
\[
  \phi(u_1,\ldots,u_D)
  =
  G_1(u_1)G_2(u_2)\cdots G_D(u_D),
  \label{eq:p14-functional-tt}
\]
\[
  G_k(u_k)=\sum_{a=1}^{n_k} C_{k,a}b_{k,a}(u_k),
  \qquad
  C_{k,a}\in\R^{r_{k-1}\times r_k},\quad r_0=r_D=1.
  \label{eq:p14-core-expansion}
\]
The storage is roughly
\[
  \sum_{k=1}^D r_{k-1}n_k r_k,
  \label{eq:p14-tt-storage}
\]
instead of
\[
  \prod_{k=1}^D n_k.
  \label{eq:p14-full-storage}
\]

\paragraph{Why this matters.}
For a chemistry or physics reader, the useful analogy is limited coupling.  If
a high-dimensional density can be described by local or moderately long-range
interactions after a good coordinate ordering, the TT ranks may stay moderate.
The TT is then not magic compression; it is a controlled low-rank description
of how coordinate groups interact.

\paragraph{Failure mode.}
If every coordinate is strongly coupled to many distant coordinates, the ranks
\(r_k\) can grow until \eqref{eq:p14-tt-storage} is no longer small.  The TT
representation remains mathematically valid, but it may stop being efficient.

\section{From Square-Root Approximation To A Filter}

The filtering density \(q_t\) is nonnegative.  Therefore its square root is
well-defined wherever \(q_t\) is finite:
\[
  s_t(u_t;\alpha)=\sqrt{q_t(u_t;\alpha)}.
  \label{eq:p14-root-target}
\]
Approximate \(s_t\) by a TT:
\[
  s_t(u_t;\alpha)\approx
  \phi_t(u_t;\alpha)
  =
  G_{t,1}(u_{t,1};\alpha)\cdots G_{t,D}(u_{t,D};\alpha).
  \label{eq:p14-root-tt}
\]
Then define
\[
  \widehat q_t(u_t;\alpha)
  =
  \phi_t(u_t;\alpha)^2+\tau_t(\alpha)\lambda_t(u_t),
  \qquad
  \lambda_t\ge0,\quad \int\lambda_t=1.
  \label{eq:p14-qhat}
\]
The normalizer and normalized joint density are
\[
  \widehat z_t(\alpha)
  =
  \int\widehat q_t(u_t;\alpha)\,\dd u_t,
  \qquad
  \widehat\pi_t(u_t;\alpha)
  =
  \frac{\widehat q_t(u_t;\alpha)}{\widehat z_t(\alpha)}.
  \label{eq:p14-zhat-pihat}
\]
The approximate filter is the marginal
\[
  \widehat p_t(x_t,\vartheta;\alpha)
  =
  \int
  \widehat\pi_t(x_t,\vartheta,x_{t-1};\alpha)
  \,\dd x_{t-1}.
  \label{eq:p14-pt-from-pihat}
\]

To compute the square part of \(\widehat z_t\), use the basis mass matrices
\[
  M_{t,k}[a,b]
  =
  \int b_{t,k,a}(u_k)b_{t,k,b}(u_k)\,\dd\nu_{t,k}(u_k).
  \label{eq:p14-mass-matrix}
\]
Starting from \(R_{t,D}=1\), define
\[
  R_{t,k-1}
  =
  \sum_{a,b=1}^{n_k}
  M_{t,k}[a,b]\,
  C_{t,k,a}R_{t,k}C_{t,k,b}^{\top},
  \qquad k=D,\ldots,1.
  \label{eq:p14-mass-contraction}
\]
Then
\[
  R_{t,0}
  =
  \int \phi_t(u_t;\alpha)^2\,\dd u_t,
  \qquad
  \widehat z_t=R_{t,0}+\tau_t
  \label{eq:p14-z-from-R}
\]
for fixed normalized \(\lambda_t\).

\paragraph{Why this matters.}
The square makes the approximation nonnegative.  The mass contraction makes
the integral computable without enumerating the full tensor grid.  The
normalization and marginalization then turn the approximation into a filter.

\paragraph{Failure mode.}
Signed or unconstrained low-rank density approximations can become negative.
The squared-TT construction avoids that problem, but it can still be inaccurate
if the square-root fit \(\phi_t\approx\sqrt q_t\) is poor.

\section{The Zhao--Cui Sequential Construction}

The sequential construction can be written as the following recursion.  Given
\(\widehat p_{t-1}\), form
\[
  q_t(x_t,\vartheta,x_{t-1})
  =
  \widehat p_{t-1}(x_{t-1},\vartheta)
  f_\alpha(x_t\mid x_{t-1},\vartheta)
  g_\alpha(y_t\mid x_t,\vartheta).
  \label{eq:p14-zc-q}
\]
Fit a square-root TT
\[
  \phi_t\approx\sqrt{q_t}.
  \label{eq:p14-zc-root-fit}
\]
Define
\[
  \widehat q_t=\phi_t^2+\tau_t\lambda_t,\qquad
  \widehat z_t=\int\widehat q_t,\qquad
  \widehat\pi_t=\widehat q_t/\widehat z_t.
  \label{eq:p14-zc-normalize}
\]
Propagate the filter by
\[
  \widehat p_t(x_t,\vartheta)
  =
  \int \widehat\pi_t(x_t,\vartheta,x_{t-1})\,\dd x_{t-1}.
  \label{eq:p14-zc-propagate}
\]

The same normalized joint density also defines conditional distributions.  If
the coordinates of \(u_t\) are ordered as \(u_1,\ldots,u_D\), then
\[
  \widehat\pi_t(u_1,\ldots,u_D)
  =
  \widehat\pi_t(u_1)
  \widehat\pi_t(u_2\mid u_1)\cdots
  \widehat\pi_t(u_D\mid u_1,\ldots,u_{D-1}).
  \label{eq:p14-conditional-factorization}
\]
A conditional Knothe--Rosenblatt map samples these coordinates by inverting
successive conditional cumulative distribution functions:
\[
  r_k
  =
  F_k(u_k\mid u_1,\ldots,u_{k-1}),
  \qquad
  u_k
  =
  F_k^{-1}(r_k\mid u_1,\ldots,u_{k-1}),
  \quad r_k\in(0,1).
  \label{eq:p14-kr}
\]
For the filtering and gradient results in this note, the map is secondary.
What matters is that the same squared-TT object supplies marginals,
conditionals, normalizers, and samples.

Zhao and Cui use this sequential squared-TT construction for state-space
learning, including the nonseparable density over
\((x_t,\theta,x_{t-1})\), the squared density
\(\phi^2+\tau\lambda\), and the marginal/conditional constructions
\cite{zhao-cui-2024}.  Cui and Dolgov provide the squared inverse Rosenblatt
transport machinery used to interpret the conditional maps
\cite{cui-dolgov-2022}.

\paragraph{Why this matters.}
The source papers do not merely provide a compression trick.  They provide a
way to use a nonnegative low-rank density representation inside the same
normalization and marginalization operations that define Bayesian filtering.

\paragraph{Failure mode.}
The conditional maps inherit the quality of the squared-TT approximation.  A
bad density approximation yields bad marginal and conditional distributions
even if the algebraic conditioning procedure is internally consistent.

\section{Why Adaptive TT Filtering Is Not Automatically A Gradient Method}

Let \(B\) denote all numerical branch choices:
\[
  B=
  \{\text{ranks, pivots, interpolation points, samples, truncations,
  factorization signs, shift rule}\}.
  \label{eq:p14-branch}
\]
For a fixed branch \(B\), the algorithm computes a scalar
\[
  \widehat\ell_T(\alpha;B)
  =
  \sum_{t=1}^T
  \{\log\widehat z_t(\alpha;B)-c_t(\alpha;B)\}.
  \label{eq:p14-branch-scalar}
\]
An adaptive code chooses a branch from the parameter value:
\[
  B=B(\alpha).
  \label{eq:p14-adaptive-branch}
\]
Thus the adaptive value path is
\[
  L_{\rm adap}(\alpha)
  =
  \widehat\ell_T(\alpha;B(\alpha)).
  \label{eq:p14-adaptive-scalar}
\]
A fixed-branch derivative computes
\[
  \left.
  \frac{\partial}{\partial\alpha_i}
  \widehat\ell_T(\alpha;B_0)
  \right|_{\alpha=\alpha_0},
  \qquad B_0=B(\alpha_0),
  \label{eq:p14-fixed-branch-derivative}
\]
not the derivative of \eqref{eq:p14-adaptive-scalar} across branch changes.

\paragraph{Why this matters.}
Adaptation is useful for filtering because the representation can respond to
the target density.  A gradient, however, must differentiate one scalar formula.
When the formula changes because \(B(\alpha)\) changes, the derivative of the
old branch is only a local branch derivative.

\paragraph{Failure mode.}
If finite-difference evaluations at \(\alpha+h e_i\) and \(\alpha-h e_i\)
choose different ranks or interpolation points, then the finite difference is
testing the adaptive value path, while the fixed-branch formula differentiates
\(\widehat\ell_T(\alpha;B_0)\).  These are different mathematical objects.

\section{The Fixed-Branch Filtering Algorithm}

A fixed branch is specified before differentiation.  For each time \(t\), fix
\[
  B_t=
  (\text{ordering},\text{domain maps},\text{basis},\text{ranks},
  \text{fitting points},\text{solver path},\lambda_t,\tau_t,c_t,
  \text{factorization signs}).
  \label{eq:p14-fixed-branch-spec}
\]
The forward step is:
\[
  q_t=\widehat p_{t-1}^{\TT}f_\alpha g_\alpha,
  \qquad
  s_t=\exp(c_t/2)\sqrt{q_t},
  \qquad
  \phi_t\approx s_t.
  \label{eq:p14-fixed-forward-target}
\]
At fixed fitting points \(u^{(j)}\), this means
\[
  s_j(\alpha)
  =
  \exp(c_t(\alpha)/2)
  \left[
  \widehat p_{t-1}^{\TT}(x_{t-1}^{(j)},\vartheta^{(j)};\alpha)
  f_\alpha(x_t^{(j)}\mid x_{t-1}^{(j)},\vartheta^{(j)})
  g_\alpha(y_t\mid x_t^{(j)},\vartheta^{(j)})
  \right]^{1/2}.
  \label{eq:p14-target-values}
\]

One possible core update is interpolation:
\[
  A_{t,k}(\alpha)g_{t,k}(\alpha)=b_{t,k}(\alpha).
  \label{eq:p14-interpolation-system}
\]
Another is fixed weighted least squares:
\[
  g_{t,k}(\alpha)
  =
  \arg\min_g
  \|W_{t,k}^{1/2}(A_{t,k}(\alpha)g-b_{t,k}(\alpha))\|_2^2.
  \label{eq:p14-ls-system}
\]
The branch must choose the value path.  The derivative later differentiates
that path, not an unspecified mixture of paths.

After the cores are obtained, compute
\[
  R_{t,0}=\int\phi_t^2,\qquad
  \widehat z_t=R_{t,0}+\tau_t,\qquad
  \widehat\ell_t^{\TT}=\log\widehat z_t-c_t.
  \label{eq:p14-fixed-forward-scalar}
\]
The normalized joint and filter are
\[
  \widehat\pi_t
  =
  \frac{\phi_t^2+\tau_t\lambda_t}{\widehat z_t},
  \qquad
  \widehat p_t^{\TT}(x_t,\vartheta)
  =
  \int\widehat\pi_t(x_t,\vartheta,x_{t-1})\,\dd x_{t-1}.
  \label{eq:p14-fixed-filter}
\]
Store the unnormalized marginal numerator
\[
  a_t(x_t,\vartheta)
  =
  \int
  \{\phi_t(x_t,\vartheta,x_{t-1})^2
  +\tau_t\lambda_t(x_t,\vartheta,x_{t-1})\}
  \,\dd x_{t-1},
  \label{eq:p14-marginal-numerator}
\]
so that
\[
  \widehat p_t^{\TT}(x_t,\vartheta)=a_t(x_t,\vartheta)/\widehat z_t.
  \label{eq:p14-filter-numerator-ratio}
\]

\paragraph{Why this matters.}
The saved numerator \(a_t\) is not bookkeeping clutter.  It is needed because
the next target \(q_{t+1}\) contains \(\widehat p_t^{\TT}\), and therefore the
gradient at \(t+1\) needs the derivative of the previous approximate filter.

\paragraph{Failure mode.}
If the implementation saves only \(\widehat\ell_t^{\TT}\) and discards the
branch objects that produced it, then a later derivative implementation cannot
prove that it differentiates the same scalar.

\section{Proposition 1: Normalized Approximate Filtering}

\begin{assumption}[Fixed-branch filtering assumptions]
For every \(t\), \(\widehat p_{t-1}^{\TT}\) is nonnegative and integrates to
one.  The transition and likelihood densities are nonnegative and measurable.
The fixed branch produces a measurable \(\phi_t\).  The defensive density
satisfies \(\lambda_t\ge0\) and \(\int\lambda_t=1\).  The defensive mass
satisfies \(\tau_t\ge0\).  Finally,
\[
  0<
  \widehat z_t
  =
  \int\{\phi_t(u)^2+\tau_t\lambda_t(u)\}\,\dd u
  <\infty.
  \label{eq:p14-pos-finite}
\]
\end{assumption}

\begin{proposition}[Fixed-branch squared-TT filtering is normalized]
Under the fixed-branch filtering assumptions, define
\[
  \widehat q_t(u)=\phi_t(u)^2+\tau_t\lambda_t(u),
  \qquad
  \widehat\pi_t(u)=\widehat q_t(u)/\widehat z_t,
  \label{eq:p14-prop1-joint}
\]
and
\[
  \widehat p_t^{\TT}(x_t,\vartheta)
  =
  \int
  \widehat\pi_t(x_t,\vartheta,x_{t-1})\,\dd x_{t-1}.
  \label{eq:p14-prop1-marginal}
\]
Then \(\widehat p_t^{\TT}\) is a nonnegative normalized density.  Starting
from a normalized prior, the recursion defines normalized approximate filters
for all times for which \eqref{eq:p14-pos-finite} holds.
\end{proposition}

\begin{proof}
Pointwise,
\[
  \widehat q_t(u)=\phi_t(u)^2+\tau_t\lambda_t(u)\ge0.
\]
Because \(0<\widehat z_t<\infty\),
\[
  \widehat\pi_t(u)\ge0,\qquad
  \int\widehat\pi_t(u)\,\dd u
  =
  \frac{\int\widehat q_t(u)\,\dd u}{\widehat z_t}
  =
  1.
\]
Using Tonelli's theorem for the nonnegative integrand,
\[
\begin{aligned}
  \int \widehat p_t^{\TT}(x_t,\vartheta)\,\dd x_t\,\dd\vartheta
  &=
  \int\int
  \widehat\pi_t(x_t,\vartheta,x_{t-1})
  \,\dd x_{t-1}\,\dd x_t\,\dd\vartheta \\
  &=
  \int \widehat\pi_t(u)\,\dd u
  =
  1.
\end{aligned}
\]
Thus the marginal filter is nonnegative and normalized.  The induction over
time is immediate: a normalized prior gives the base case, and the displayed
argument advances the property from \(t-1\) to \(t\).
\end{proof}

\paragraph{Why this matters.}
The proposition proves that the approximation is a proper probability density
after every completed step.  It does not prove that the density is accurate;
that depends on the approximation quality of \(\phi_t\).

\paragraph{Failure mode.}
If \(\widehat z_t=0\), \(\widehat z_t=\infty\), or the numerical branch creates
a nonmeasurable or nonintegrable object, the normalization argument fails.

\section{Proposition 2 In Layers: Same-Scalar Gradient}

This section expands the gradient proof in layers.  Each layer contributes one
piece of the final derivative.

\subsection*{Layer A: derivative of the declared scalar}

The fixed-branch scalar is
\[
  \widehat\ell_T^{\TT}(\alpha)
  =
  \sum_{t=1}^T
  \{\log\widehat z_t(\alpha)-c_t(\alpha)\}.
  \label{eq:p14-declared-scalar}
\]
If \(\widehat z_t(\alpha)>0\), then
\[
  \partial_i
  \{\log\widehat z_t(\alpha)-c_t(\alpha)\}
  =
  \frac{\partial_i\widehat z_t(\alpha)}
       {\widehat z_t(\alpha)}
  -
  \partial_i c_t(\alpha).
  \label{eq:p14-layerA}
\]

\paragraph{Why this matters.}
The gradient problem reduces to computing \(\partial_i\widehat z_t\) and
\(\partial_i c_t\) for the same branch that produced \(\widehat z_t\) and
\(c_t\).

\subsection*{Layer B: derivative of the squared-TT normalizer}

For fixed normalized \(\lambda_t\),
\[
  \widehat z_t
  =
  \int\{\phi_t(u;\alpha)^2+\tau_t(\alpha)\lambda_t(u)\}\,\dd u.
  \label{eq:p14-z-layerB}
\]
Assuming differentiation may pass through the integral,
\[
\begin{aligned}
  \partial_i\widehat z_t
  &=
  \int
  2\phi_t(u;\alpha)\partial_i\phi_t(u;\alpha)\,\dd u
  +
  \partial_i\tau_t(\alpha)\int\lambda_t(u)\,\dd u \\
  &=
  2\int
  \phi_t(u;\alpha)\partial_i\phi_t(u;\alpha)\,\dd u
  +
  \partial_i\tau_t(\alpha).
\end{aligned}
  \label{eq:p14-layerB}
\]
If \(\lambda_t\) depends on \(\alpha\), the extra term
\[
  \tau_t\int\partial_i\lambda_t(u;\alpha)\,\dd u
  \label{eq:p14-lambda-extra}
\]
must be included.

\paragraph{Why this matters.}
The normalizer derivative depends on the square-root TT derivative, not on a
separate derivative of a normalized filter.  This is why the derivative path
must track \(\partial_i\phi_t\).

\subsection*{Layer C: derivative of the TT square root}

The square-root TT is
\[
  \phi_t(u)
  =
  G_{t,1}(u_1)G_{t,2}(u_2)\cdots G_{t,D}(u_D).
  \label{eq:p14-layerC-tt}
\]
The product rule gives
\[
  \partial_i\phi_t(u)
  =
  \sum_{k=1}^D
  G_{t,1}(u_1)\cdots
  \partial_iG_{t,k}(u_k)\cdots
  G_{t,D}(u_D).
  \label{eq:p14-layerC-product}
\]
With
\[
  G_{t,k}(u_k)
  =
  \sum_{a=1}^{n_k} C_{t,k,a}b_{t,k,a}(u_k),
  \label{eq:p14-layerC-core}
\]
one has, for fixed basis functions,
\[
  \partial_iG_{t,k}(u_k)
  =
  \sum_{a=1}^{n_k}
  (\partial_iC_{t,k,a})b_{t,k,a}(u_k).
  \label{eq:p14-layerC-coredot}
\]

\paragraph{Why this matters.}
The derivative of the high-dimensional function is obtained by differentiating
one core at a time.  This is the same low-rank coupling idea as Section 3, now
applied to sensitivities.

\subsection*{Layer D: derivative of the mass contraction}

The forward mass contraction is
\[
  R_{t,D}=1,\qquad
  R_{t,k-1}
  =
  \sum_{a,b}
  M_{t,k}[a,b]C_{t,k,a}R_{t,k}C_{t,k,b}^{\top}.
  \label{eq:p14-layerD-R}
\]
Assuming fixed mass matrices, the differentiated contraction is
\[
\begin{aligned}
  \partial_iR_{t,k-1}
  =
  \sum_{a,b}M_{t,k}[a,b]\Big[
  &(\partial_iC_{t,k,a})R_{t,k}C_{t,k,b}^{\top}
  +C_{t,k,a}(\partial_iR_{t,k})C_{t,k,b}^{\top}\\
  &+C_{t,k,a}R_{t,k}(\partial_iC_{t,k,b})^{\top}
  \Big],
\end{aligned}
  \label{eq:p14-layerD-Rdot}
\]
with \(\partial_iR_{t,D}=0\).

\paragraph{Why this matters.}
The scalar \(R_{t,0}\) used in the forward likelihood is differentiated by the
same backward contraction.  This prevents the gradient from silently switching
to another normalizer.

\subsection*{Layer E: derivative of the core solve}

For an interpolation branch,
\[
  A_{t,k}g_{t,k}=b_{t,k}.
  \label{eq:p14-layerE-interp}
\]
Differentiating gives
\[
  A_{t,k}\partial_i g_{t,k}
  =
  \partial_i b_{t,k}
  -
  (\partial_iA_{t,k})g_{t,k}.
  \label{eq:p14-layerE-interpdot}
\]
For a weighted least-squares branch,
\[
  g_{t,k}
  =
  \arg\min_g
  \|W_{t,k}^{1/2}(A_{t,k}g-b_{t,k})\|_2^2,
  \label{eq:p14-layerE-ls}
\]
the normal equations are
\[
  N_{t,k}g_{t,k}=A_{t,k}^{\top}W_{t,k}b_{t,k},
  \qquad
  N_{t,k}=A_{t,k}^{\top}W_{t,k}A_{t,k}.
  \label{eq:p14-layerE-normal}
\]
With fixed \(W_{t,k}\),
\[
\begin{aligned}
  N_{t,k}\partial_i g_{t,k}
  =
  &(\partial_i A_{t,k})^\top W_{t,k}
  (b_{t,k}-A_{t,k}g_{t,k})\\
  &+
  A_{t,k}^{\top}W_{t,k}
  (\partial_i b_{t,k}-(\partial_iA_{t,k})g_{t,k}).
\end{aligned}
  \label{eq:p14-layerE-lsdot}
\]

\paragraph{Why this matters.}
The derivative of a core is not guessed.  It is obtained by differentiating
the same linear algebra problem that produced the core.

\subsection*{Layer F: derivative of the previous filter}

The target at time \(t\) contains the previous filter:
\[
  q_t
  =
  \widehat p_{t-1}^{\TT}f_\alpha g_\alpha.
  \label{eq:p14-layerF-q}
\]
Where the factors are positive,
\[
  \partial_i\log q_t
  =
  \partial_i\log\widehat p_{t-1}^{\TT}
  +
  \partial_i\log f_\alpha
  +
  \partial_i\log g_\alpha.
  \label{eq:p14-layerF-logq}
\]
Using the stored numerator from \eqref{eq:p14-marginal-numerator},
\[
  \widehat p_{t-1}^{\TT}(v;\alpha)
  =
  \frac{a_{t-1}(v;\alpha)}{\widehat z_{t-1}(\alpha)},
  \label{eq:p14-layerF-ratio}
\]
so
\[
  \partial_i\log\widehat p_{t-1}^{\TT}(v;\alpha)
  =
  \frac{\partial_i a_{t-1}(v;\alpha)}{a_{t-1}(v;\alpha)}
  -
  \frac{\partial_i\widehat z_{t-1}(\alpha)}
       {\widehat z_{t-1}(\alpha)}.
  \label{eq:p14-layerF-prev}
\]
The derivative \(\partial_i a_{t-1}\) is another marginal contraction, with
one previous TT core differentiated at a time.

\paragraph{Why this matters.}
The gradient is recursive because the filter is recursive.  Time \(t\) cannot
be differentiated correctly unless the derivative of the time \(t-1\) filter
is carried forward.

\begin{assumption}[Fixed-branch differentiability]
The fixed-branch filtering assumptions hold.  The branch choices are fixed.
The model densities, previous-filter numerator, coordinate maps, basis maps,
core coefficients, defensive mass, and stabilizing constant are differentiable
in \(\alpha_i\).  The displayed differentiations may pass through the finite
contractions and integrals.  The fixed interpolation or least-squares systems
are nonsingular, or a differentiable pseudoinverse branch is declared.
\end{assumption}

\begin{proposition}[Same-scalar fixed-branch gradient]
Under the fixed-branch differentiability assumptions, the derivative of
\[
  \widehat\ell_T^{\TT}(\alpha)
  =
  \sum_{t=1}^T
  \{\log\widehat z_t(\alpha)-c_t(\alpha)\}
\]
is
\[
  \partial_i\widehat\ell_T^{\TT}
  =
  \sum_{t=1}^T
  \left[
  \frac{\partial_iR_{t,0}+\partial_i\tau_t}
       {R_{t,0}+\tau_t}
  -
  \partial_i c_t
  \right],
  \label{eq:p14-final-score}
\]
for the direct mass-contraction branch with fixed normalized
\(\lambda_t\).  If \(\lambda_t\) is parameter-dependent, the additional term
in \eqref{eq:p14-lambda-extra} must be included.
\end{proposition}

\begin{proof}
Layer A proves the derivative of each log-normalizer contribution.  Layer B
expresses \(\partial_i\widehat z_t\) in terms of
\(\phi_t\), \(\partial_i\phi_t\), and \(\partial_i\tau_t\).  Layer C gives
\(\partial_i\phi_t\) by differentiating one TT core at a time.  Layer D shows
that the derivative of the square mass \(R_{t,0}\) is computed by the same
mass-contraction recursion used in the forward pass, with product-rule terms.
Layer E gives the core sensitivities for the fixed interpolation or fixed
least-squares branch.  Layer F supplies the recursive derivative of the
previous approximate filter inside \(q_t\).  Substituting
\(\widehat z_t=R_{t,0}+\tau_t\) into Layer A gives
\eqref{eq:p14-final-score}.
\end{proof}

\paragraph{Failure mode.}
If the rank, pivot set, fitting points, branch signs, or least-squares system
change with \(\alpha\), the derivative above remains the derivative of the
frozen local branch, not of the adaptive code path.

\section{Algorithmic Gradient Recipe}

\paragraph{Inputs.}
For each \(t\), provide
\[
  B_t,\quad \widehat p_{t-1}^{\TT},\quad f_\alpha,\quad g_\alpha,\quad
  \partial_i\log f_\alpha,\quad \partial_i\log g_\alpha,\quad
  y_t.
\]
The branch data \(B_t\) includes bases, ranks, fitting points, mass matrices,
linear-system type, defensive density, defensive mass, shift rule, and
factorization signs.

\paragraph{Forward pass.}
For \(t=1,\ldots,T\):
\[
  q_t=\widehat p_{t-1}^{\TT}f_\alpha g_\alpha,
  \qquad
  s_t=\exp(c_t/2)\sqrt{q_t}.
\]
Solve either \eqref{eq:p14-interpolation-system} or
\eqref{eq:p14-ls-system} for the TT cores.  Compute
\[
  R_{t,0},\quad
  \widehat z_t=R_{t,0}+\tau_t,\quad
  \widehat\ell_t^{\TT}=\log\widehat z_t-c_t,
\]
and store \(a_t\) from \eqref{eq:p14-marginal-numerator}.

\paragraph{Sensitivity pass.}
For each parameter coordinate \(\alpha_i\) and time \(t\):
\[
  \partial_i\log q_t
  =
  \partial_i\log\widehat p_{t-1}^{\TT}
  +
  \partial_i\log f_\alpha
  +
  \partial_i\log g_\alpha.
\]
Differentiate the chosen core solve, producing \(\partial_iC_{t,k,a}\).
Propagate \(\partial_iR_{t,k}\) backward with
\eqref{eq:p14-layerD-Rdot}.  Form
\[
  \partial_i\widehat z_t
  =
  \partial_iR_{t,0}+\partial_i\tau_t
\]
and
\[
  \partial_i\widehat\ell_t^{\TT}
  =
  \frac{\partial_i\widehat z_t}{\widehat z_t}
  -
  \partial_i c_t.
  \label{eq:p14-step-score}
\]
Carry the filter sensitivity forward through
\[
  \partial_i(\widehat q_t/\widehat z_t)
  =
  \frac{\partial_i\widehat q_t}{\widehat z_t}
  -
  \frac{\widehat q_t\,\partial_i\widehat z_t}{\widehat z_t^2},
  \label{eq:p14-normalized-density-dot}
\]
then marginalize over \(x_{t-1}\).

\paragraph{Finite-difference parity.}
For a frozen branch \(B_0\), check
\[
  \frac{
  \widehat\ell_T^{\TT}(\alpha+h e_i;B_0)
  -
  \widehat\ell_T^{\TT}(\alpha-h e_i;B_0)}
  {2h}
  \approx
  \partial_i\widehat\ell_T^{\TT}(\alpha;B_0).
  \label{eq:p14-parity}
\]
If the branch is not frozen, this test no longer checks the same derivative.

\paragraph{Failure mode.}
The most dangerous implementation error is a value-gradient mismatch: the
forward pass computes one scalar, while the derivative pass differentiates a
slightly different scalar.  The parity test is designed to catch that error.

\section{What The Construction Gives And What Must Be Tested}

Mathematically, the construction gives
\[
  \widehat q_t=\phi_t^2+\tau_t\lambda_t,\qquad
  \widehat z_t=\int\widehat q_t,\qquad
  \widehat p_t^{\TT}=\int\widehat q_t/\widehat z_t\,\dd x_{t-1}.
\]
Proposition 1 proves that this is a normalized approximate filter.  Proposition
2 proves that, on a fixed branch,
\[
  \nabla_\alpha
  \sum_t\{\log\widehat z_t-c_t\}
\]
is computed by differentiating the same branch.

What must still be tested is numerical:
\[
  \|\widehat q_t-q_t\|,\qquad
  \max_k r_{t,k},\qquad
  \kappa(A_{t,k})\ \text{or}\ \kappa(N_{t,k}),\qquad
  \text{finite-difference parity error}.
\]
These quantities answer different questions.  Approximation error tests
whether the filter is accurate.  Rank growth tests whether the TT idea is
efficient.  Conditioning tests whether the linear algebra is stable.  Parity
tests whether the implemented gradient differentiates the same scalar.

\paragraph{Why this matters.}
The derivation makes Zhao--Cui-style TT filtering a credible candidate because
it defines a normalized approximate filter and a fixed-branch gradient.  It
does not by itself show that the candidate wins on the target model.  That
requires numerical experiments and implementation tests.

\appendix
\section{Source And Code Anchors}

The derivations above are written in BayesFilter notation.  The source anchors
below record where the corresponding objects appear in the inspected sources.

\begin{itemize}[leftmargin=*]
\item Zhao--Cui equations (1)--(3): state-space model notation.
\item Zhao--Cui equations (9)--(12) and Algorithm 1: sequential density over
  \((x_t,\theta,x_{t-1})\), TT approximation, and integration architecture.
\item Zhao--Cui equation (13) and nearby Lemma 1: squared-TT density
  \(\phi^2+\tau\lambda\).
\item Zhao--Cui Algorithm 2: sequential use of squared-TT approximations.
\item Zhao--Cui Propositions 2 and 4: marginal and conditional constructions.
\item Zhao--Cui Section 4.1: normalizer/evidence discussion.
\item Cui--Dolgov Sections 2--3: squared inverse Rosenblatt and functional TT
  transport substrate.
\item Inspected MATLAB snapshot: the scalar path adds
  \(\log(\texttt{sirt.z})-\texttt{const}\), and the TT marginalization stores a
  normalizer of the form square-mass plus defensive mass.  These code
  references show consistency of scalar structure; the mathematical
  justification is the derivation in the main text.
\end{itemize}

\begin{thebibliography}{9}
\bibitem{zhao-cui-2024}
Zhao, Y. and Cui, T. (2024).
\newblock Tensor-train methods for sequential state and parameter learning in
state-space models.
\newblock \emph{Journal of Machine Learning Research}, 25(244):1--51.

\bibitem{cui-dolgov-2022}
Cui, T. and Dolgov, S. (2022).
\newblock Deep composition of tensor-trains using squared inverse Rosenblatt
transports.
\newblock \emph{Foundations of Computational Mathematics}, 22:1863--1922.
\newblock DOI: 10.1007/s10208-021-09537-5.
\end{thebibliography}

\end{document}
